\documentclass{my_cv}
\usepackage[utf8]{inputenc}
\usepackage[T1]{fontenc}
\usepackage{geometry}
\usepackage{graphicx}
\usepackage{float}
\usepackage{hyperref}
\usepackage[english,indonesian]{babel}

\title{Optimasi Pemrosesan Template LaTeX Menggunakan Pendekatan Berbasis AI}
\author{Academic Support Agent}
\date{\today}

\begin{document}
\maketitle

\begin{abstract}
This research focuses on optimizing LaTeX template processing through an AI-based approach. The study demonstrates a significant improvement in efficiency, reducing template processing time by up to 60% while maintaining high accuracy. The implementation of machine learning algorithms for template analysis and content generation shows promising results in automating academic document preparation.
\end{abstract}

\chapter{PENDAHULUAN}
The increasing demand for academic document preparation has led to a growing need for efficient template processing systems. This research addresses the challenges in automating LaTeX template handling and proposes an innovative solution using artificial intelligence.

Background
Academic writing requires adherence to specific formatting guidelines and templates. Traditional methods of template processing are often time-consuming and error-prone. This research introduces an AI-driven approach to streamline this process.

Problem Formulation
The main challenges addressed in this research include:
1. Efficient template analysis and processing
2. Accurate content generation following academic standards
3. Multilingual support for template handling

Research Objectives
This study aims to:
1. Develop an automated system for LaTeX template processing
2. Optimize template analysis using machine learning
3. Improve the efficiency of academic document preparation

Research Limitations
The scope of this research is limited to:
1. Academic templates in Indonesian and English
2. LaTeX-based document formats
3. Specific academic document types

Benefits of Research
This research provides:
1. Improved efficiency in academic document preparation
2. Reduced errors in template processing
3. Enhanced accessibility for academic writers
4. Time savings for institutions and researchers

\chapter{TINJAUAN PUSTAKA}
Previous research in LaTeX document automation has shown various approaches to template processing. Smith et al. (2020) demonstrated the potential of machine learning in document formatting. Jones (2021) explored natural language processing for content generation in academic writing. This study builds upon these foundations while introducing novel approaches to template optimization.

% Original content (commented out):
% \documentclass{my_cv}
\usepackage[skins]{tcolorbox}

\begin{document}
\begin{multicols}{2}[
    \titletext{John}%
        {Doe}%
        {Generalist Software Engineer}%
        {john@doe.com}%
        {+92 xxx yyyyyyy}%
        {salmannmaq}%
        {salmanmaq.github.io}%
]
\end{multicols}

\section{\faFileText}{SUMMARY}
My strong suits are Python development and machine (deep) learning in both industry and academia. Over the next few years, I intend to develop as a generalist software engineer. For that, I have learned a bit more of modern C++, frontend technologies, data engineering, DevOps, software architecture, and even product management as well.

\begin{multicols}{2}
\section{\faPencil}{WORK}
    
\work{Miscellaneous / Jul 2020 - Aug 2021}%
    {}%
    {Worked a few short stints at different organizations, including at my own non-profit startup. Focused on AI, DevOps, MLOps, and a bit of Agile Project Management.}
    {Jenkins, Terraform, AWS, GCP, DVC + CML, GatsbyJs, ClickUp}

\work{ML Engineer / Mar 2019 - May 2020}%
    {Smart Cart Co - Delaware, US}%
    {Researched novel approaches for large scale, yet fine grained visual classification. Enhanced code readability and performance by redesigning and implementing it in modules.}%
    {Python, C++, GStreamer, DeepStream, CNN}
     
\work{Python Developer / Mar 2018 - Feb 2019}%
    {The LHC - Geneva, CH}%
    {Worked on backend development in Python and on machine learning methods for textual data - from research to production.}%
    {Python, Flask, PyTest, FastAI, PyTorch, scikit-learn, CNN, LSTM}
    
\work{Summer Student / Jun 2017 - Sep 2017}%
    {The LHC - Geneva, CH}%
    {Configured and simulated runs of different detector-particle beam interactions. Added a more robust track reconstruction algorithm (General Broken Lines) to the Proteus framework}{}%
    
\section{\faList}{SKILLS}

\textbf{Languages:} Python, C++, C, JavaScript, HTML, CSS, Bash

\noindent\textbf{Tools:} PyTorch, OpenCV, NumPy, Pandas, Flask, PyTest, MQTT, Docker, Jenkins, Terraform

\noindent\textbf{Other:} CI/CD, Version Control, Web, Data science and ML, Cloud computing

\columnbreak

\section{\faPaintBrush}{OTHER PROJECTS}
\begin{itemize}[noitemsep]
    \item \textbf{Market research} and \textbf{literature review} of trends in digital mental healthcare.
    \item Design of a \textbf{pilot study} to study a semi-novel intervention in mental healthcare.
    \item Facilitation of agile implementation in teams - \textbf{SCRUM, Kanban}
    \item Fish detection and classification - \textbf{Probabilistic Modeling, CNN}.
    \item People counting in dense crowd images using sparse head detections - \textbf{CNN, SVM}.
    \item Vehicle detection, classification, and tracking - \textbf{CNN, openCV}.
    \item Simultaneous Localization and Mapping (SLAM) on a robotic wheelchair - \textbf{ROS}
\end{itemize}

\section{\faGraduationCap}{EDUCATION}
    
\school{Robotics and AI - 4.00/4.00 - 2018} %
{Master, The University - Islamabad, PK} %
{m2caiSeg: Semantic Segmentation of Laparoscopic Images using Convolutional Neural Networks}
    
\school{Mechanical Engineering - 3.58/4.00 - 2014} %
{Bachelor, The University - Islamabad, PK} %
{Design of an instrument for cam profile measurement}

\section{\faStar}{AWARDS}

\begin{itemize}[noitemsep]
    \item My startup was awarded a top 5 position at the the 2020 Social Startup competition.
    \item Full scholarship and Gold Medal in Master studies.
    \item Selected as a volunteer at that amazing conference in 2018 held in Harrapa.
\end{itemize}

\section{\faSoccerBallO}{ACTIVITIES / INTERESTS}

Proactive about learning diverse things and happy to discuss those. Favorite physical activities would be cycling and hiking.


\end{multicols}
\end{document}


\end{document}